% !TEX TS-program = XeLaTeX
% !TEX encoding = UTF-8 Unicode

\chapter{相关研究与技术基础}
\label{chap:related}
\defaultfont

\section{WiFi 人体感知研究}
WiFi 人体感知利用无线信号与人体之间的相互作用进行行为、姿态和身份识别。人体进入无线传播区域后，会改变信号传播路径，使接收端观测到的信道响应发生变化。已有研究表明，CSI 能够刻画子载波级别的信道变化，在室内定位、动作识别和手势识别等任务中具有较好的表现\cite{halperin2011csi,indotrack2017,signfi2018}。

与基于图像的姿态检测方法相比，WiFi 感知不直接采集人体图像，具有更好的隐私保护特性；与可穿戴设备相比，WiFi 感知不需要用户佩戴额外传感器，适合室内无感化部署。因此，将 WiFi CSI 引入脊柱姿态检测任务，有助于构建低成本、非接触的姿态感知方法。

\section{CSI 信道状态信息}
CSI 描述无线信号在发射端与接收端之间传播时的信道响应。对于正交频分复用系统，不同子载波上的 CSI 能够反映不同频率分量受到的衰落、反射和多径影响。CSI 通常以复数形式表示，其中幅度反映信号强弱，相位反映信号传播过程中的相对偏移。

人体姿态变化会引起信号传播路径变化，使 CSI 序列在时间维度上呈现出不同模式。脊柱姿态相关的身体形态变化主要体现在躯干轮廓和运动方式差异上，这些差异会改变人体对无线信号的遮挡和反射方式，从而形成可被模型学习的 CSI 特征。

\section{CSI 相位校正}
CSI 原始相位通常受到采样时钟偏移、载波频率偏移、硬件延迟和环境噪声影响。若直接使用原始相位进行分类，模型容易受到设备状态和环境变化干扰。相位校正的目标是削弱与人体姿态无关的公共相位误差，增强由人体姿态变化引起的相对相位差异。

基于 CSI 比值的相位校正方法通过不同天线链路之间的复数比值构造相对相位特征。本文采用接收天线之间的相位比作为主要特征：以一根接收天线为参考，将其他接收天线的复数 CSI 与参考天线 CSI 相除，再取复数比值的幅角作为相位比。该方法保留链路之间的相对变化信息，并在一定程度上抵消共同相位偏移，适合用于多天线 WiFi 感知任务。

\section{传统机器学习分类模型}
CSI 姿态检测任务中的样本数量通常有限，而单个样本可以从多天线、多子载波和时间维度构造出较高维度的特征。对于此类小样本、高维、非线性分类问题，传统机器学习模型仍具有较强适用性。支持向量机（Support Vector Machine，SVM）能够通过核函数处理非线性分类边界\cite{svm1995}；Random Forest 通过集成多棵决策树降低单棵树的方差\cite{randomforest2001}；ExtraTrees 在随机森林基础上进一步增强特征选择和切分阈值的随机性，从而构造差异更大的树模型并提升集成泛化能力\cite{extratrees2006}。

本文最终采用 ExtraTrees 作为脊柱姿态检测分类器。该模型能够处理高维统计特征，不依赖严格的特征缩放过程，并能通过多棵随机树的集成投票表达复杂的非线性特征组合。对于当前 750 个 CSI 文件构成的数据集，ExtraTrees 在测试集上取得了较好的分类表现。

\section{深度学习时间序列模型}
深度学习模型能够从原始或半结构化特征中自动学习分类表示。对于 CSI 姿态检测任务，输入数据具有明显的时间序列属性，因此可使用一维卷积网络、循环神经网络或时间卷积网络进行建模。时间卷积网络（Temporal Convolutional Network，TCN）通过卷积核在时间维度上滑动提取局部动态特征，并通过多层结构扩大感受野\cite{tcn2018}。

本文实验过程中曾使用 TCN 作为窗口级时间序列分类基线。实验结果显示，在当前数据规模下，TCN 容易学习到受试者和采集环境中的局部细节，而测试集表现低于基于文件级统计特征的树集成模型。因此，本文将 TCN 作为模型对比的一部分，并将相位比统计特征与 ExtraTrees 分类器确定为最终方案。
